\begin{theorem}[一阶收缩常数对对称序贯判别律的闭包]\label{thm:conclusion-golden-sprt-delta-closure}
记 \(W\) 为黄金互补二元信道，并定义
\[
\delta:=\delta(W),
\qquad
\rho_{\max}:=\rho_{\mathrm m}(W),
\qquad
\eta_{\mathrm{SDPI}}:=\eta_I(W).
\]
则
\[
\delta=\rho_{\max}=\varphi^{-3},
\qquad
\eta_{\mathrm{SDPI}}=\varphi^{-6}=\delta^2.
\]
进一步，对任意阈值 \(T\ge 1\)，结论 \ref{thm:conclusion-golden-sprt-risk-time-elimination} 中的对称序贯判别律可完全改写为
\[
\varepsilon_T=\frac{1}{1+\delta^{-T/3}},
\]
\[
\mathbb E_{\theta_+}[\tau_T]
=
\delta^{-1}T\cdot
\frac{\delta^{-T/3}-1}{\delta^{-T/3}+1}.
\]
因此，同一条一阶收缩常数 \(\delta\) 同时决定单步 TV/相关收缩、互信息 SDPI 常数，以及多步对称 SPRT 的整条风险--样本量曲线。
\end{theorem}

\begin{proof}
由定理 \ref{thm:pom-golden-binary-channel-contraction}，
\[
\delta(W)=\rho_{\mathrm m}(W)=\varphi^{-3},
\qquad
\eta_I(W)=\varphi^{-6}=\delta(W)^2.
\]
将
\[
\varphi^T=(\varphi^{-3})^{-T/3}=\delta^{-T/3},
\qquad
\varphi^3=\delta^{-1}
\]
代入定理 \ref{thm:conclusion-golden-sprt-risk-time-elimination} 中的闭式
\[
\varepsilon_T=\frac{1}{1+\varphi^T},
\qquad
\mathbb E_{\theta_+}[\tau_T]
=
\varphi^3T\cdot\frac{\varphi^T-1}{\varphi^T+1},
\]
便得到所示重写公式。证毕。
\end{proof}

\endinput
